\documentclass[12pt]{article}
\usepackage[utf8]{inputenc}
\usepackage[T1]{fontenc}
\usepackage{hyperref}
\title{Fe'fe' Language Resources for K-12 Curriculum Development}
\author{Compiled by Claude Sonnet 5.5 (claude-sonnet-5-5) for Andr\'as Kornai}
\date{}
\begin{document}
\maketitle

\section{Introduction}
Fe'fe' (Fe'efe'e, Nufi, Bafang) is a Bamileke language of the Eastern Grassfields branch, spoken by about 140,000 people in the Haut-Nkam department, West Region, Cameroon. Missionaries began writing it in the 1920s, and the current orthography follows the General Alphabet of Cameroonian Languages (1978). No comprehensive modern dictionary or school curriculum was found.

\section{Primary Educational Resources}
No textbooks or primers were identified in the search results. The search results do state that the language was at one time taught in schools, but no materials were located.

\section{Linguistic and Reference Materials}
\begin{itemize}
\item Fe'fe' language overview (Wikipedia): \url{https://en.wikipedia.org/wiki/Fe%CA%BCfe%CA%BC_language}
\item Omniglot page on Fe'fe' writing: \url{https://omniglot.com/writing/fefe.htm}
\item Nka'kayii, Nzambu' 1973, French grammar of Bamileke-Fe'efe'e (Glottolog reference record): \url{https://glottolog.org/resource/reference/id/147572}
\item Tchamna and Deukam, grammar of Nfi (Fe'efe'e), in French (library catalogue record): \url{https://catalog.pascolibraries.org/Hoopla/15423698}
\end{itemize}

\section{Religious and Community Resources}
Christian texts were translated by French missionaries from the 1920s according to the Wikipedia article; no specific current edition with a URL was found.

\section{Digital and Audio Resources}
None found.

\section{Curriculum Development Context}
The language is stable with a large community and is taught as a subject in some schools. A curriculum would build on the AGLC orthography and the grammars above, and would need new materials, ideally developed with local Fe'fe' speakers and teachers.
\end{document}